% This is text associated with the open textbook "Open Electromagnetics"
% See immediately below for author, date
% This work is licensed under the Creative Commons Attribution-ShareAlike 4.0 International License. (CC BY SA 4.0) 
% To view a copy of the license, visit http://creativecommons.org/licenses/by-sa/4.0/

%ORB TITLE Wave Power in a Lossless Medium
%ORB AUTHOR 0        # S.W. Ellingson, Virginia Tech (ellingson@vt.edu)
%ORB DATE 20180807   # yyyymmdd
%ORB ABSTRACT Plane Wave Power Density in a Lossless Medium 

\label{m00041_Plane_Wave_Power_Density_in_a_Lossless_Medium} % <-- Must match name of this file and module directory name.  Don't mess with this!
\index{waves!power density|(}
\index{power density|see{waves}}

In many applications involving electromagnetic waves, one is less concerned with the instantaneous values of the electric and magnetic fields than the \emph{power} associated with the wave. 
In this section, we address the issue of how much power is conveyed by an electromagnetic wave in a lossless medium.
The relevant concepts are readily demonstrated in the context of uniform plane waves, 
\index{waves!plane}
as shown in this section.
A review of Section~\ref{m0039_Uniform_Plane_Waves_Characteristics} (``Uniform Plane Waves: Characteristics'') is recommended before reading further.

Consider the following uniform plane wave, described in terms of the phasor representation of its electric field intensity:
%
\begin{equation}
\widetilde{\bf E} = \hat{\bf x}E_0 e^{-j\beta z}
\label{m0041_eE}
\end{equation}
%
Here $E_0$ is a complex-valued constant associated with the source of the wave, and $\beta$ is the positive real-valued propagation constant.  
Therefore, the wave is propagating in the $+\hat{\bf z}$ direction in lossless media.

The first thing that should be apparent is that the amount of power conveyed by this wave is infinite.
The reason is as follows. If the power passing through any finite area is greater than zero, then the total power must be infinite because, for a uniform plane wave, the electric and magnetic field intensities are constant over a plane of infinite area.
In practice, we never encounter this situation because all practical plane waves are only ``locally planar'' (see Section~\ref{m0142_Types_of_Waves} for a refresher on this idea).
Nevertheless, we seek some way to express the power associated with such waves.

The solution is not to seek total power, but rather power per unit area.
This quantity is known as the \emph{spatial power density}, or simply ``\emph{power density},'' and has units of W/m$^2$.\footnote{Be careful: The quantities \emph{power spectral density} (W/Hz) and \emph{power flux density} (W/(m$^2\cdot$Hz)) are also sometimes referred to as ``power density.''  In this section, we will limit the scope to \emph{spatial} power density (W/m$^2$).} 
Then, if we are interested in total power passing through some finite area, then we may simply integrate the power density over this area.
Let's skip to the answer, and then consider where this answer comes from.
It turns out that the \emph{instantaneous} power density of a uniform plane wave is the magnitude of the \emph{Poynting vector}
\index{Poynting vector}
%
\begin{equation}
{\bf S} \triangleq {\bf E} \times {\bf H}
\label{m0041_eS}
\end{equation}
%
Note that this equation is dimensionally correct; i.e. the units of ${\bf E}$ (V/m) times the units of ${\bf H}$ (A/m) yield the units of spatial power density (V$\cdot$A/m$^2$, which is W/m$^2$).
Also, the direction of ${\bf E} \times {\bf H}$ is in the direction of propagation (reminder: Section~\ref{m0039_Uniform_Plane_Waves_Characteristics}), which is the direction in which the power is expected to flow.
Thus, we have some compelling evidence that $\left|{\bf S}\right|$ is the power density we seek.
However, this is not proof -- for that, we require the \emph{Poynting Theorem}, 
\index{Poynting theorem}
which is a bit outside the scope of the present section, but is addressed in the ``Additional Reading'' at the end of this section.

A bit later we're going to need to know ${\bf S}$ for a uniform plane wave, so let's work that out now.
From the plane wave relationships (Section~\ref{m0039_Uniform_Plane_Waves_Characteristics}) we find that the magnetic field intensity associated with the electric field in Equation~\ref{m0041_eE} is
%
\begin{equation}
\widetilde{\bf H} = \hat{\bf y}\frac{E_0}{\eta} e^{-j\beta z}
\label{m0041_eH}
\end{equation}
%
where $\eta=\sqrt{\mu/\epsilon}$ is the real-valued impedance of the medium.
Let $\psi$ be the phase of $E_0$; i.e., $E_0 = |E_0|e^{j\psi}$.
Then
%
\begin{flalign}
{\bf E} &=  \mbox{Re}\left\{ \widetilde{\bf E} e^{j\omega t} \right\}  \nonumber \\
        &= \hat{\bf x}\left|E_0\right| \cos\left( \omega t - \beta z + \psi \right)
\end{flalign}
%
and
%
\begin{flalign}
{\bf H} &=  \mbox{Re}\left\{ \widetilde{\bf H} e^{j\omega t} \right\}  \nonumber \\
        &= \hat{\bf y}\frac{\left|E_0\right|}{\eta} \cos\left( \omega t - \beta z + \psi \right)
\end{flalign}
%
Now applying Equation~\ref{m0041_eS},
%
\begin{equation}
{\bf S} =\hat{\bf z} \frac{|E_0|^2}{\eta} \cos^2\left( \omega t - \beta z + \psi \right)
\end{equation} 

As noted earlier, $\left|{\bf S}\right|$ is only the instantaneous power density, which is still not quite what we are looking for.
What we are actually looking for is the \emph{time-average} power density $S_{ave}$ -- that is, the average value of $\left|{\bf S}\right|$ over one period $T$ of the wave.
This may be calculated as follows:
%
\begin{flalign}
S_{ave} &= \frac{1}{T} \int_{t=t_0}^{t_0+T} |{\bf S}| dt \nonumber \\
        &= \frac{|E_0|^2}{\eta} \frac{1}{T} \int_{t=t_0}^{t_0+T}  \cos^2\left( \omega t - ks + \psi \right)  dt
\end{flalign}
%
Since $\omega =2\pi f = 2\pi/T$, the definite integral equals $T/2$.  We obtain
%
\begin{equation}
\boxed{ S_{ave} = \frac{|E_0|^2}{2\eta} }
\label{m0041_eWEPWD}
\end{equation}
%
It is useful to check units again at this point.  Note (V/m)$^2$ divided by $\Omega$ is W/m$^2$, as expected.  
%
%ORB HIGHLIGHT
\begin{mdframed}[backgroundcolor=yellow!20,nobreak=true] 
Equation~\ref{m0041_eWEPWD} is the time-average power density (units of W/m$^2$) associated with a sinusoidally-varying uniform plane wave in lossless media.
\end{mdframed}

Note that Equation~\ref{m0041_eWEPWD} is analogous to a well-known result from electric circuit theory. Recall the time-average power $P_{ave}$ (units of W) associated with a voltage phasor $\widetilde{V}$ across a resistance $R$ is 
%
\begin{equation}
P_{ave} = \frac{|\widetilde{V}|^2}{2R}
\label{m0041_ePWLM-EA}
\end{equation}
%
which closely resembles Equation~\ref{m0041_eWEPWD}.
The result is also analogous to the result for a voltage wave on a transmission line (Section~\ref{m0090_Power_Flow_on_Transmission_Lines}), for which: 
%
\begin{equation}
P_{ave} = \frac{|V_0^+|^2}{2Z_0}
\end{equation}
%
where $V_0^+$ is a complex-valued constant representing the magnitude and phase of the voltage wave, and $Z_0$ is the characteristic impedance 
\index{characteristic impedance}
of the transmission line.
\index{waves!transmission line analogy}

Here is a good point at which to identify a common pitfall. $|E_0|$ and $|\widetilde{V}|$ are the {\it peak magnitudes} of the associated real-valued physical quantities.  However, these quantities are also routinely given as {\it root mean square} (``rms'') quantities.  Peak magnitudes are greater by a factor of $\sqrt{2}$, so Equation~\ref{m0041_eWEPWD} expressed in terms of the rms quantity lacks the factor of $1/2$.  

%%%%%%%%%%%%%%%%%%%%%%%%%%%%%%%%%%%%%%%%%%%%%%%
%ORB EXAMPLE
~\\ \begin{example} 
\begin{mdframed}[backgroundcolor=brown!20,linecolor=brown!20]\raggedright
Power density of a typical radio wave.
\\~\\
A radio wave transmitted from a distant location may be perceived locally as a uniform plane wave if there is no nearby structure to scatter the wave; a good example of this is the wave arriving at the user of a cellular telephone in a rural area with no significant terrain scattering.  The range of possible signal strengths varies widely, but a typical value of the electric field intensity arriving at the user's location is 10~$\mu$V/m rms.  What is the corresponding power density?

{\bf Solution}:  From the problem statement, $|E_0|=10~\mu$V/m rms.
We assume propagation occurs in air, which is indistinguishable from free space at cellular frequencies.  
If we use Equation~\ref{m0041_eWEPWD}, then we must first convert $|E_0|$ from rms to peak magnitude, which is done by multiplying by $\sqrt{2}$.  Thus:
%
\begin{flalign}
S_{ave} &= \frac{|E_0|^2}{2\eta} \cong \frac{\left(\sqrt{2} \cdot 10 \times 10^{-6}~\mbox{V/m}\right)^2}{2 \cdot 377~\Omega} \nonumber \\
        &\cong 2.65 \times 10^{-13}~\mbox{W/m$^2$} \nonumber
\end{flalign}
%
Alternatively, we can just use a version of Equation~\ref{m0041_eWEPWD}, which is appropriate for rms units:
%
\begin{flalign}
S_{ave} &= \frac{|E_{0,rms}|^2}{\eta} \cong \frac{\left(10 \times 10^{-6}~\mbox{V/m}\right)^2}{ 377~\Omega} \nonumber \\
        &\cong 2.65 \times 10^{-13}~\mbox{W/m$^2$} \nonumber
\end{flalign}
%
Either way, we obtain the correct answer, 0.265~pW/m$^2$ (that's {\it picowatts} per square meter).

\end{mdframed}
% note figures will need to go between "\end{mdframed}" and "\end{example}"
\end{example}
%%%%%%%%%%%%%%%%%%%%%%%%%%%%%%%%%%%%%%%%%%%%%%%


Considering the prevalence of phasor representation, it is useful to have an alternative form of the Poynting vector which yields time-average power by operating directly on field quantities in phasor form.  This is ${\bf S}_{ave}$, defined as:
%
\begin{equation}
{\bf S}_{ave} \triangleq \frac{1}{2} \mbox{Re} \left\{ \widetilde{\bf E} \times \widetilde{\bf H}^* \right\}
\end{equation}
%
(Note that the magnetic field intensity phasor is conjugated.)  
The above expression gives the expected result for a uniform plane wave.
Using Equations~\ref{m0041_eE} and \ref{m0041_eH}, we find
%
\begin{equation}
{\bf S}_{ave} = \frac{1}{2} \mbox{Re} \left\{
  \left( \hat{\bf x}E_0 e^{-j\beta z} \right)
\times
  \left( \hat{\bf y}\frac{E_0}{\eta} e^{-j\beta z} \right)^* \right\}
\end{equation}
%
which yields
%
\begin{equation}
{\bf S}_{ave} = \hat{\bf z}\frac{|E_0|^2}{2\eta} 
\end{equation}
%
as expected.  
%The details of the derivation are not difficult and are left as an exercise for the reader.

%Although we have developed the Poynting vector concept from a scenario involving a single uniform plane wave, it should be noted that the concept applies to {\it any} wave or combination of waves.  This is verifiably so since in the scenario we considered we did not need to specify the volume of $\mathcal{V}$ nor the area of $A$ or $B$.  Thus, these dimensions could be made arbitrarily small, such that any wave appears planar and uniform to arbitrarily good accuracy.  

{\bf Additional Reading:}
\begin{itemize}
\item ``\href{https://en.wikipedia.org/wiki/Poynting_vector}{Poynting vector}'' on Wikipedia.
\item ``\href{https://en.wikipedia.org/wiki/Poynting's_theorem}{Poynting's Theorem}'' on Wikipedia.
\end{itemize}    

\index{waves!power density|)}
